\section{Distributional sources: compact returns and spacetime translations}
\label{sec:distributional}
We now test two further, explicitly specified source relations.
The first stays in the original finite-slab state and retains its
radial electric potential. The second explicitly assumes an
infinite-volume continuum vacuum. Neither changes the scope of
the preceding exclusions into a prohibition on all YM sources.

\subsection{Short-support mixed identities and compact returns}
For nonzero $h\in\cD$, put
\begin{equation}\label{eq:shrinking-test}
h_R(x)=R^{1/2}h(Rx),\qquad f_R=R^{-2}\T h_R\in\cD^0,
\qquad c_h=\norm{h''}_2^2>0.
\end{equation}
These tests shrink in support; they are not the modulations
of Proposition~\ref{prop:highfreq}. Their exact normalizations are
\begin{equation}
\norm{f_R}_2^2=c_h+\frac{\norm{h'}_2^2}{2R^2}
 +\frac{\norm h_2^2}{16R^4},\qquad
\widehat f_R(\tau)=R^{-5/2}(\tau^2+1/4)\widehat h(\tau/R).
\end{equation}
\begin{lemma}[Shrinking-support mixed test]\label{lem:shrinking}
With the complete contact convention \eqref{eq:contact-growth},
\begin{align}
Q[f_R]&=c_h\log R+d_h+o(1),&
d_h&=\frac1{2\pi}\int u^4\log\frac{|u|}{2\pi}
                         |\widehat h(u)|^2\dd u,\label{eq:shrinking-diag}\\
Q(f_R,U_tf_R)&=-b(t)\norm{f_R}_2^2+O_{h,t}(R^{-5})
\quad(t\ne0),\label{eq:shrinking-off}\\
Q(f_R,g)&=O_{h,g}(R^{-5/2})\quad(g\in\cD^0\ \hbox{fixed}).
\label{eq:shrinking-weak}
\end{align}
Here $b(t)=\Lambda(a)/\sqrt a$ if $|t|=\log a$ for an integer
$a\ge2$, and is zero otherwise.
\end{lemma}
\begin{proof}
The diagonal prime sum vanishes once the support diameter is
less than $\log2$. Scale the archimedean integral to obtain
\[
\frac1{2\pi}\int m_+(Ru)(u^2+(4R^2)^{-1})^2
                              |\widehat h(u)|^2\dd u.
\]
Use \eqref{eq:contact-growth} for $|u|\ge R^{-1}$ and boundedness
of $m_+$ on the remaining interval. The $u^4\log|u|$ term is
integrable, the other tails are Schwartz, and the polynomial
correction is $O(R^{-2}\log R)$. This proves
\eqref{eq:shrinking-diag}. Adding a constant c to the multiplier
would change $d_h$ by $cc_h$; the contact has not been reset.

At fixed nonzero t the supports are eventually separated.
Move both copies of $\T$ onto the smooth off-diagonal kernel
$-n_\Gamma(|x-y|)$. Its derivatives are bounded there, and
$R^{-4}\norm{h_R}_1^2=R^{-5}\norm h_1^2$.
Local discreteness of the prime shifts leaves only the shift
exactly cancelling t, giving \eqref{eq:shrinking-off}, including
when t is a prime-power separation. For fixed g, the Fourier
formula and its rapid decay prove \eqref{eq:shrinking-weak}
for the archimedean term. There are only finitely many prime
terms on the enclosing supports; moving $\T$ onto each
translate of g gives the same estimate.
\end{proof}

\begin{theorem}[No compact return for an exact source]\label{thm:compact-return}
Suppose a continuous law $J:\cD^0\to\cH$ into a positive
physical Hilbert space has pairing Q. On
$\mathcal S=\overline{J\cD^0}$ let $V_t$ implement
$JU_t=V_tJ$, either as the induced action or as a restriction
of a specified physical unitary group. Then, for every $L\ne0$,
\begin{equation}\label{eq:essential-circle}
\sigma_{\rm ess}(V_L|_{\mathcal S})=\{z:|z|=1\}.
\end{equation}
In particular no nonzero Laurent polynomial in $V_L$ is
compact on $\mathcal S$. Thus $V_L-e^{\ii\phi}I$ cannot be
compact there or on an ambient sector containing it.
\end{theorem}
\begin{proof}
Set $x_R=Jf_R/\sqrt{\log R}$. Lemma~\ref{lem:shrinking} and
density give $x_R\rightharpoonup0$ and
$\norm{x_R}^2\to c_h$. For every nonzero integer k,
$\ip{x_R}{V_L^kx_R}\to0$ by \eqref{eq:shrinking-off}.
Consequently, for $p(z)=\sum_jp_jz^j$,
\begin{equation}
\norm{p(V_L)x_R}^2\longrightarrow c_h\sum_j|p_j|^2.
\end{equation}
This contradicts compactness for every nonzero p. Uniform
polynomial approximation extends the formula to continuous
F on the circle, with limit
$c_h\int|F(e^{\ii\theta})|^2\dd\theta/(2\pi)$.
If an open arc missed the essential spectrum, a nonzero
continuous F supported there would give compact $F(V_L)$
by continuous functional calculus modulo compact operators,
again a contradiction.
\end{proof}
The proof uses the geometric form directly, not RH or a zero
spectrum. Its weakly null family is a contradiction test,
not a proposed physical limit retaining norms. Pure point
spectrum alone is not excluded: Kronecker flows can have
the full-circle essential spectrum required here. Compact
resolvent of a generator alone does not imply a compact return.

\subsection{Candidate A: uncompensated electric waves in both polarizations}
Keep the fixed state of Section~\ref{sec:boundary}. For
$P=\cos\theta I+\ii\sin\theta\,\omega\cdot\sigma$, set
$\Sigma(P)=\omega\cdot\sigma$, ignoring the null endpoints.
The existing scalar and adjoint reference-color channels give
the actual closed sector
\begin{equation}
\cH_{\rm pol}=\{F_0(\theta)I+F_1(\theta)\Sigma(P)\}
\simeq L^2(\mu_\rho)\oplus L^2(\mu_\rho),\qquad
\dd\mu_\rho=\tfrac2\pi\rho\sin^2\theta\dd\theta.
\end{equation}
The trace/2 norm makes the channels orthogonal. Compress the
actual electric form divided by $\ell=4$ to this sector:
\begin{equation}\label{eq:polar-electric}
q_\rho[F]=\frac2\pi\int_0^\pi\rho\{
\sin^2\theta(|F_0'|^2+|F_1'|^2)+2|F_1|^2\}\dd\theta.
\end{equation}
Each distinct link supplies the same SU(2) Dirichlet term;
the angular factor is $\sum_j|\nabla_{S^2}\omega_j|^2=2$.
Take the closed forms and their Friedrichs operators
$E_0,E_1\ge0$. This is a form compression, not a claim that
the full $E_\nu$ reduces the sector.

With $v=\sqrt\rho\sin\theta$ and $WF=\sqrt{2/\pi}\,vF$,
\begin{equation}\label{eq:polar-liouville}
WE_0W^{-1}=-\partial^2+v''/v,\qquad
WE_1W^{-1}=-\partial^2+2\csc^2\theta+v''/v.
\end{equation}
The first has Dirichlet endpoints and the second the
Friedrichs inverse-square endpoints. The retained potential
$v''/v$ is smooth and bounded. No compensation is performed.
The candidate translation is the electric-wave prescription
\begin{equation}\label{eq:electric-wave}
G_{\rm el}=\sqrt{1+E_0}\oplus(-\sqrt{1+E_1}),\qquad
V_x=e^{-\ii xG_{\rm el}}.
\end{equation}
Both frequency directions are supplied inside the existing
physical sector. The unit shift uses the fixed Casimir units;
any fixed positive shift or common speed rescaling gives the
same conclusion below. Identifying arithmetic x with this
electric-wave parameter is the candidate's extra source relation,
not a consequence of YM dynamics or physical Euclidean time.

Let $e_{s,n}$ be its actual electric eigenvectors with positive
frequencies $\omega_{s,n}$, and let a distributional seed have
coefficients satisfying
\begin{equation}
\sum_{s,n}(1+\omega_{s,n}^2)^{-r}|b_{s,n}|^2<\infty
\end{equation}
for some finite r. Evaluation at a fixed interior holonomy
angle, for example $\pi/2$, and its angular derivatives are
independent geometric examples, by Sobolev evaluation.
With $\epsilon_0=1$, $\epsilon_1=-1$, define
\begin{align}
J_bf&=\sum_{s,n}b_{s,n}\widehat f(\epsilon_s\omega_{s,n})e_{s,n},\\
\ip{J_bf}{J_bg}_\nu
&=\sum_{s,n}|b_{s,n}|^2
\overline{\widehat f(\epsilon_s\omega_{s,n})}
\widehat g(\epsilon_s\omega_{s,n}).\label{eq:electric-wave-pair}
\end{align}
Fourier decay proves strong convergence of spectral cutoffs,
smooth-test continuity and the displayed positive pairing.
The eigenvalues and coefficients of a fixed geometric seed
retain the actual state dependence. Finite-order derivatives
are included by increasing r. There is no subtraction or
renormalization, and no winding phase is discarded: this
candidate uses electric eigenvectors rather than scaling
packets. Any admissible winding-modified seed is included
with all its coefficients and obeys the same exclusion.

\begin{proposition}[Failure of the radial electric-wave candidate]
No source covariant under \eqref{eq:electric-wave} has the
complete Weil pairing, even without a finite-order seed bound.
\end{proposition}
\begin{proof}
Put $V_\rho=1+v''/v$, a bounded real potential.
The Haar comparison operators for $1+E_s$ have eigenvalues
$n^2$, starting at $n=1$ for the scalar channel and $n=2$
for the adjoint channel. For the latter use
$A=\partial-\cot\theta$, with
$A^*A=-\partial^2-1$ and
$AA^*=-\partial^2+2\csc^2\theta-1$.
The vectors $A\sin(n\theta)$, $n\ge2$, are complete;
the missing adjoint kernel $1/\sin\theta$ is not in $L^2$.
The min-max principle for the bounded perturbation $V_\rho$
therefore gives
\begin{equation}
|\omega_{s,n}^2-n^2|\le\norm{V_\rho}_\infty,\qquad
\omega_{s,n}=n+O_\rho(n^{-1}).
\end{equation}
Thus $V_{2\pi}-I$ is compact on both channels. Apply
Theorem~\ref{thm:compact-return}.
\end{proof}
This pure-point candidate avoids the earlier absolutely
continuous current obstruction and requires no winding
intertwiner. It fails because of its asymptotic spacing.
The full many-link electric operator need not have this
spacing; no conclusion about that operator follows here.

\subsection{Candidate B: translations in a continuum vacuum}
We now explicitly change to continuum YM on Minkowski
$\R^{1,3}$. Assume continuum existence and full OS
reconstruction with a separable physical Hilbert space,
a strongly continuous positive-energy unitary Poincar\'e
representation, and a unique translation-invariant vacuum.
A physical mass gap is allowed but unnecessary below.
For a fixed nonzero real four-vector a the candidate relation is
\begin{equation}\label{eq:spacetime-source}
JU_t=e^{-\ii tP\cdot a}J.
\end{equation}
It describes translating a fixed physical preparation along a
straight spacetime direction. For $b=A\Omega_{\rm vac}$ and
a polynomial p, derivative smearing has actual pairing
\begin{equation}
J_bf=\{p(\ii\lambda)\widehat f(\lambda)\}(P\cdot a)b,
\qquad
\ip{J_bf}{J_bg}=\int\overline{\widehat f}\widehat g
                         |p(\ii\lambda)|^2\dd\mu_b^a(\lambda).
\end{equation}
The combined multiplier is bounded on each test and gives
strong spectral-cutoff convergence. More singular seeds
are permitted if smearing gives actual Hilbert vectors.

\begin{lemma}[Spacetime-line spectral type]\label{lem:spacetime-ac}
Under these continuum representation assumptions,
$P\cdot a$ has absolutely continuous spectrum off the
vacuum for every $a\ne0$.
\end{lemma}
\begin{proof}
Let E be the joint momentum spectral measure and choose a
finite control measure
$\mu=\sum_n2^{-n}\ip{e_n}{E(\cdot)e_n}$ from an orthonormal
basis. It has the same null sets as E, and covariance makes
its measure class Lorentz invariant. Average it over a smooth
compactly supported nonnegative Lorentz-group density positive
near the identity. The average has the same null sets:
one direction follows from covariance, and the other from
continuity and nonnegativity of
$\ip{U(\Lambda)e_n}{E(B)U(\Lambda)e_n}$ at the identity.

For every nonzero forward-cone p, the real-analytic function
$\Lambda\mapsto a\cdot\Lambda p$ is nonconstant. Its critical
set has Haar measure zero. The submersion/coarea theorem
therefore makes the inverse image of each Lebesgue-null set
Haar-null. Tonelli and equivalence of the averaged measure
prove the assertion for all $p\ne0$. At zero momentum,
uniqueness leaves only the vacuum.
\end{proof}
This proof includes spacelike and lightlike a and requires
neither a particle decomposition nor asymptotic completeness.
It uses infinite-volume Lorentz covariance, which the fixed
slab does not possess. Mass-shell representation theory gives
the same interpretation \cite{bekaert-boulanger}.

Consequently each vector coefficient is a vacuum constant
plus a function in $C_0(\R)$. If it equalled $C_*$, boundedness
would imply RH by Theorem~\ref{thm:probe-criterion}.
Under that consequence \eqref{eq:C-positive} has zero Ces\`aro
mean and positive mean square. The first removes the vacuum
constant; the second contradicts $C_0$ decay. Hence no exact
source can satisfy \eqref{eq:spacetime-source}, even for
distributional preparations whose smeared vectors exist.
This verifies a physical spectral hypothesis rather than
assuming clustering in an unidentified arithmetic variable.
Neither candidate establishes the missing non-geometric
source relations or compact arithmetic feasibility.
